
% --- SLIDE 3: MLFLOW COMPONENTS ---
\begin{frame}{MLflow: Four Components}

    \centering
    \renewcommand{\arraystretch}{1.35}
    {\scriptsize
    \begin{tabular}{p{2.6cm}|p{9.0cm}}
        \toprule
        \textbf{Component} & \textbf{What it is} \\
        \midrule
        \textbf{Tracking}   & An API and UI for logging parameters, code versions, metrics and artifacts
                              when you run ML code, and for visualising the results afterwards. \\
        \textbf{Projects}   & A standard format for packaging and sharing reproducible data-science code,
                              with its entry points and environment declared, so someone else can rerun it. \\
        \textbf{Models}     & A standard packaging format with ``flavours''. Anything supporting the
                              \texttt{python\_function} flavour can be served as a REST endpoint, batch-scored
                              in Spark, or exported to a cloud serving platform. \\
        \textbf{Model Registry} & A centralised model store with versioning, aliases, tags and
                              annotations, plus lineage back to the run that produced each version. \\
        \bottomrule
    \end{tabular}}

    \vspace{0.24em}
    \begin{itemize}
        \scriptsize
        \item They are designed to work together but each one is usable alone. \textbf{Start with Tracking}
              --- it is three lines of code and it is where all the value is at first.
        \item MLflow is open source and runs entirely on your laptop: no account, no network, no cost.
    \end{itemize}

    \vspace{0.09em}
    {\scriptsize Source: MLflow documentation, \href{https://mlflow.org/docs/latest/ml/}{mlflow.org/docs/latest/ml}.
    See also Zaharia et al., ``Accelerating the Machine Learning Lifecycle with MLflow,''
    \textit{IEEE Data Engineering Bulletin} 41(4), 2018.}

\end{frame}
